The four classic types of threats~\cite{wohlin2012experimentation} are discussed below, each with the steps taken to reduce it.

\textbf{Construct validity.} The main construct, the agent behind an issue, comes from a heuristic attribution that its authors describe as a candidate with false positives by design~\cite{ghaleb2026codage}. It was kept as given, as the study brief requires, but the study reports for every agent how its issues are linked to it, repeats the key tests on the high-precision subset (S3) and after removing \NPreRelease{} issues created before the agent's public release (S1), and the researcher validated a stratified random sample of 99 issues by hand (nine per tested agent, rated blind). In this sample, \ValCorrectPct\% of the attributions are correct (95\% CI \ValCorrectCI), ranging from \ValCorrectAiderK{} of \ValCorrectAiderN{} for Aider to all nine for Claude Code, Devin and Kiro. Most errors (\ValWrongOtherK{} of \ValWrongOtherN) come from text signals, for example the word ``Raider'' matched as Aider, or the ordinary word ``sweep'' matched as the agent Sweep, which is the weakness the dataset's documentation warns about. The own-bot flag, on which the closure gap in RQ1 rests, was right for all \ValOwnBotPrecisionK{} issues it flagged, but it caught only \ValOwnBotRecallK{} of the \ValOwnBotRecallN{} sampled issues that the agent itself had written (\ValOwnBotRecallPct\%, 95\% CI \ValOwnBotRecallCI). It misses, for example, issues that people posted with a ``Generated with Claude Code'' footer. RQ1 therefore compares issues opened by the agent's own bot account with all of the agent's other issues, a group that also contains some agent-written and some misattributed issues. The rule was kept unchanged rather than tuned after seeing the validation. Some attributed issues are automated content that merely mentions an agent, for example \NDigest{} daily paper digests posted by \texttt{github-actions[bot]} that mention SWE-agent or Aider; the S3 subset excludes all but \NDigestInSThree{} of them. Time-to-close measures whether an issue was closed, not whether the problem was solved, and an issue closed as \emph{not planned} is not a fix; each closure reason is therefore also tested separately. The label categories come from published rules defined for this study, the language flag is a heuristic that was checked by hand on 40 issues (\ValNonEnK{} of \ValNonEnN{} flagged non-English and \ValEnK{} of \ValEnN{} flagged English issues were classified correctly), and the RQ3 topic labels are an interpretation of word lists and sample issues. They were checked by hand on five issues per topic, and only the two fine topics whose labels passed are interpreted (Section~\ref{sec:rq3}). All judgements were made by the author alone; the blind samples and the answers are in the replication package, so that others can re-rate them.

\textbf{Internal validity.} This is an observational study, so all results are associations, not effects of the agents. Bot authorship is confounded with the agent: CodeRabbit and Devin issues are mostly bot-written and the others are not. Comparisons within the same agent (RQ1, RQ2) and the human-authored subset (S2) address this; in S2, the agent effect on time-to-close shrinks from medium to small. Issues of different agents also differ in age; age was checked and does not explain the closure gap, and the survival analysis accounts for the different time each issue has been open. Many issues from one repository could inflate the evidence, but no repository holds more than \TopRepoMax\% of any tested agent's issues, and one value per agent and repository (S4) gives the same effect size.

\textbf{External validity.} The data are one extract of public GitHub issues for \NAgents{} agents, without any filter for repository popularity or activity, so personal and toy projects are included. The agents change quickly, and the results describe them as they were used in this period. Five of the tested agents have fewer than 110 issues, which limits what can be said about them individually, and RQ3 covers English issues only.

\textbf{Conclusion validity.} Time-to-close is heavily skewed, so all tests are rank-based, with an effect size reported for every test and a confidence interval with every rate and median. Families of tests are Holm-corrected, and sparse tables are tested by permutation. Statistical claims require at least 30 closed issues per agent; Selvanayagam and Ghaleb used a stricter threshold of 50~\cite{selvanayagam2026aitoai}, which would remove five of the agents studied here. Single LDA runs were not reproducible on this corpus, so the topics come from Ensemble LDA, graded by their reproducibility, and the agent--topic association is confirmed with BERTopic. The fine-view rule and the ``partly reproduced'' grade were defined after the headline result; this is stated explicitly, and the headline level uses only the rule fixed in advance. The agent--topic association becomes weaker when very short issues are added. Finally, in gensim~4.3 the setting \(\eta = \)``auto'' does not change the fitted topics, so fixed values of \(\eta\) were tuned instead. To guard against errors, every number in this report is generated from the result files, and independent scripts re-compute the key statistics.
